\section{Discussion: Research Gaps and Implications}
\label{sec:governance}

The comparison points to several gaps that none of the reviewed studies addresses on its own:
\begin{enumerate}
    \item \textbf{Generation and detection are studied separately.} Vykopal et al.\ test general detectors of machine-generated text on LLM-written disinformation~\cite{vykopal2024disinformation}, while propaganda-technique detection is evaluated on human-written news from SemEval-2020~\cite{dasanmartino2020semeval, gaeta2025towards}. None of the reviewed studies tests whether LLM-based propaganda detectors work on \emph{LLM-generated} disinformation. A direct test would combine the two settings.
    \item \textbf{Simulating variation, not only correctness.} The hyper-accuracy distortion~\cite{aher2023using} and the limited steerability of opinions~\cite{santurkar2023whose} both show that aligned models do not reproduce the spread of human answers. Studies that use LLMs as simulated respondents need methods to produce realistic variation, and must report how close the simulated distribution is to human data. The own experiment adds that spread and correctness have to be checked separately: a narrow spread can also come from a crowd that agrees on a wrong value (Section~\ref{sec:experiment}); repeating it across model sizes of one open family would show where the distortion begins.
    \item \textbf{Language and population coverage.} The simulation, opinion and generation studies use English only, and the opinion study covers only US groups. MULTITuDE shows that detection has to be evaluated across languages~\cite{macko2023multitude}; the same holds for opinion alignment and disinformation generation.
    \item \textbf{Fast-changing models.} All studies evaluate specific model versions. Because safety behaviour differs strongly between models~\cite{vykopal2024disinformation} and changes with alignment~\cite{aher2023using, santurkar2023whose}, results need to be re-checked for new models. Vykopal et al.\ suggest that LLM-based evaluation could make such repeated audits cheaper~\cite{vykopal2024disinformation}. % SOURCE: Vykopal et al. p.8 conclusion (GPT-4 can partially automate evaluation)
\end{enumerate}
For practice, the reviewed evidence suggests two cautious implications. Simulated survey or experiment participants should be treated as a complement to, not a replacement for, human data. Detection of machine-generated or manipulative text should be combined with other measures (such as source verification and platform policies), because no reviewed detector is reliable enough to be used alone.
